% =============================================================================
% Appendix subsection for freezing + encoder ablation.
% Append to SpatialWhisperer_ICML/appendix.tex as a new \section.
% Numbers from results_crc_pathocell.csv (2026-05-14).
% =============================================================================

\section{Encoder choice and freezing configuration}\label{appendix:freezing_encoder}

To understand performance of our method subject to established
hyperparameters, we evaluated distinct freezing and encoder parameters.

\paragraph{Freezing axis}
Following the LiT~\citep{Zhai2021-ao} convention, the three-letter code
$(\phi_{\mathcal{G}}, \phi_{\mathcal{T}}, \phi_{\mathcal{I}})$ specifies
locked ($L$) versus unfrozen ($U$) towers. We evaluate $LUL$ (main paper),
$LLL$ (only projection heads trained), and $LLU$ (image tower unfrozen).
We attempted $ULL$ (transcriptome unfrozen), but the configuration did
not fit in the available 80\,GB H100 memory at the batch sizes that
produced stable training, so we omit it; the smaller-batch fallback
diverged numerically.

\paragraph{Encoder axis}
We swap the bridge encoder Geneformer~\citep{Theodoris2023-nt} for an
8-layer UCE~\citep{Rosen2023-uce}.\footnote{Checkpoint
\texttt{KuanP/uce-cxg-2025-baseline-8l-512d} from
\url{https://huggingface.co/KuanP/uce-cxg-2025-baseline-8l-512d}.}
Image and text encoders are held at UNI2 and BioBERT.

\paragraph{Results}
Table~\ref{tab:freezing_encoder_appendix} reports macro AUROC on
PathoCellBench (CRC; 13 classes; presence-based aggregation of
App.~\ref{appendix:metric_discussion}).
Freezing only the projection heads ($LLL$) costs 8.3 points absolute,
confirming that text-encoder optimization is the dominant source of gain
on this task, consistent with prior reports in CellWhisperer
\citep{Schaefer2025} and LiT \citep{Zhai2021-ao}.
Unfreezing the image tower ($LLU$) costs 5.4 points absolute relative to
$LUL$: at one epoch the additional capacity does not offset the loss of
the cached UNI2 representations.
The Geneformer $\to$ UCE swap yields comparable performance to Geneformer
at this scale (1.3 points absolute below $LUL$); the bridge encoder
choice is not the bottleneck.

\begin{table}[h]
    \centering
    \small
    \caption{Hyperparameter ablation: effect of freezing configuration
    and bridge-encoder choice on cell-type prediction (PathoCellBench
    CRC; 13 classes; seed=0; one full epoch). \emph{Macro AUROC}
    averages per-class presence-based AUROC across the 13 classes.
    \emph{Macro F1} and \emph{Recall@5} are averaged across the 109
    fields-of-view. Best per column \textbf{bold}.}
    \label{tab:freezing_encoder_appendix}
    \begin{tabular}{lccc}
        \toprule
        Configuration                       & Macro AUROC      & Macro F1 & Recall@5 \\
        \midrule
        $LUL$ Geneformer (main paper)       & \textbf{0.6259}  & \textbf{0.1061} & 0.5575 \\
        $LLL$ (proj.\ heads only)           & 0.5428           & 0.0878 & 0.5350 \\
        $LLU$ (image unfrozen)              & 0.5721           & 0.0967 & \textbf{0.5788} \\
        $LUL$ UCE bridge                    & 0.6134           & 0.0615 & 0.5114 \\
        \bottomrule
    \end{tabular}
\end{table}

% =============================================================================
% Forward references to insert into main.tex
% =============================================================================
%
% main.tex line 274 (after the existing sentence ending "to the paired datasets."):
%   APPEND " App.~\ref{appendix:freezing_encoder} ablates this choice and
%           an alternative bridge encoder."
%
% main.tex line 289 (after "trained for 4 epochs."):
%   APPEND " See App.~\ref{appendix:freezing_encoder} for sensitivity of
%           this choice and an encoder-axis ablation."


% =============================================================================
% Email draft to Kuan (separate file: kuan_uce_email.txt — paste into mail
% client when ready). Asks whether the 8-layer CELLxGENE 2025 checkpoint
% is appropriate to cite in this context.
% =============================================================================
